\chapter{Open Gaps}
\label{sec:origins-gaps}
% ===================================================================

\begin{enumerate}[label=\textbf{Gap \arabic*.}, leftmargin=*, itemsep=8pt]

  \item \textbf{Density of replicating terms
        (Conjecture~\ref{conj:density}).}
        The argument that the added replication machinery can always
        be hidden from any given HML formula requires a careful
        analysis of which formulae can probe the replication channel.
        In particular, if $\phi = \langle K \rangle \psi$ where $K$
        involves the replication channel, the argument fails.
        A precise density theorem would need to show that such
        formulae form a meager (first-category) set in the logical
        topology, so that ``generically'' $\phi$ does not probe
        replication.

  \item \textbf{Defining the elementary terms.}
        Proposition~\ref{prop:assembly-index} requires a designated
        set $E$ of elementary terms.
        In the chemical setting, these are atoms or monomers.
        In the Rho calculus, the natural candidates are the
        behavioral primes of Part~II (Section~4): the irreducible
        terms from which all others are built by interaction.
        Confirming that this identification gives assembly indices
        matching those of Assembly Theory requires working through
        the combinatorics of the Rho calculus prime factorization.

  \item \textbf{The basin of attraction.}
        The claim that a population entering a neighborhood of
        $\mathrm{Rep}$ is ``sucked in'' to the replicating fixed
        point requires a Lyapunov-style argument: a function on the
        process space that is monotonically decreasing along
        trajectories near the fixed point.
        The natural candidate is the distance $d_{\HML}$ to
        $\mathrm{Rep}$, but showing it decreases requires
        assumptions on the interaction dynamics.

  \item \textbf{The compact cluster as the new geosphere.}
        The identification of Smith's fourth geosphere with the
        compact coherence cluster centered on the replicating fixed
        point requires showing that this cluster is indeed compact
        (in the sense of Part~II) and that it persists under the
        perturbations of the selection dynamics.

  \item \textbf{Crossover as a variant of comm.}
        The claim that genetic recombination is one definitional
        step from the comm rule requires specifying precisely what
        ``fragment substitution'' means in the Rho calculus and
        showing that the resulting operator has the recombination
        properties of a genetic crossover (fitness inheritance,
        schema theorem, etc.).

  \item \textbf{Gradient-coupled rules and the weight map.}
        Definition of gradient-coupled rules modifies the weight
        map of Chapter~\ref{ch:weight} by introducing an external
        parameter $\Phi$.
        Half of this gap has closed since it was written.
        Time-varying weight maps are not an extension of the
        framework but its defining feature: the update functions of
        Definition~\ref{wt:def:augmented} vary the map at every
        step, and Theorem~\ref{wt:thm:markov} says the process
        remains an exactly simulable chain because the map is part
        of the state.
        What remains open is the \emph{externally driven} case.
        An update reading a parameter the configuration does not
        carry leaves the chain, and the repair is to bring $\Phi$
        inside --- either as a component of the configuration or as
        traffic on its own channels with its own keys.
        Which of those is right is a modelling question we have not
        settled, and only the first keeps the process Markov.

  \item \textbf{Metabolic viability (Conjecture~\ref{conj:viable}).}
        The condition that catabolic inflow exceeds anabolic cost
        in expectation requires computing the path integral of a
        driven open system --- a substantially harder problem than
        the closed-system path integral of Part~I.
        The quantum-jump simulator of Chapter~\ref{ch:weight}
        provides a simulation method, but an analytic viability
        criterion requires further development.

  \item \textbf{Trophic level as causal depth
        (Proposition~\ref{prop:trophic}).}
        The identification of trophic level with
        $d^o_{\min}(K_\Phi, B)$ requires showing that the
        account-flow graph $\mathcal{F}$ is faithfully captured
        by the open synchronization tree structure.
        In particular, it requires that every account transfer
        corresponds to a minimal-context transition, which is
        plausible but needs verification.

  \item \textbf{Non-well-founded scopes.}
        Remark~\ref{sco:rmk:munu} argues that a budgeted learner can
        inhabit only a least fixed point, since membership must
        terminate.
        But an ecology with no bottom is exactly what one would model
        with the greatest fixed point, and the reflection tower is
        bounded below only by no-self-code, which is a fact about
        names rather than about ecologies.
        Is there a coherent notion of a scope that is
        non-well-founded and yet surveyable stratum by stratum?

  \item \textbf{The copy series under merger.}
        \S\ref{eng:S8} shows that merger sets $\bt = 0$ and destroys
        factorisation.
        What does it do to the copy series of
        Definition~\ref{scm:def:series}?
        The expectation is that it collapses two adjacent strata into
        one, which would make the series a direct diagnostic for
        whether a composite has merged or merely composed --- and
        would give the composition/merger distinction of \chComp an
        observable signature.

  \item \textbf{Affordable scope and the individuation fixed point.}
        A learner's region is bounded by what it can pay to survey;
        Remark~\ref{eng:rem:fixpoint} makes individuation a fixed
        point.
        Both are circular in the same way --- the region determines
        what can be afforded and what can be afforded determines the
        region.
        Do they have a common solution, and is the individual the
        least one?

  \item \textbf{The survey is an assay.}
        A survey visits channels and decides bisimilarity at a depth;
        that is an experiment, and by \S\ref{sci:S4} an experiment has
        three outcomes, the third being budget-relative.
        So a survey has a $\bot$ column, and a copy number is properly
        a triple: confirmed copies, confirmed non-copies, and channels
        the budget could not settle.
        What is the right generalisation of
        Definition~\ref{scm:def:cnn} that carries the third column,
        and does the biosignature criterion survive it?

\end{enumerate}

% ===================================================================
